%% Starts a fresh page: the table below is too tall to share one with a stranded heading.
\clearpage
\section{What does not work yet}
\label{sec:short:limits}

A document that only said the first five sections would be a pitch. This one is the reason to
believe the rest. Each row is a finding from the full paper, stated once and without softening; the
last column is what, if anything, is being done about it.

\begin{center}\small
\renewcommand{\arraystretch}{1.25}
\begin{tabular}{@{}>{\raggedright\arraybackslash}p{3.15cm}p{4.85cm}p{5.2cm}@{}}
\toprule
& The limit & Why it matters, and the response \\
\midrule
\textbf{Hosting pays nothing extra} &
PNR pays every eligible node in a tier the same share, regardless of what it runs. The marginal
reward for hosting an application is exactly zero. &
Free-riding is individually rational; what compels hosting is the eligibility gate, and that gate is
currently bound to neither the operator nor the machine. This is deliberate --- paying by placement
would make the claim race a race for money --- and the answer is attested execution, which is why
that work is a precondition rather than a follow-on. \\
\textbf{Four parties reach the removal threshold} &
At the measured collateral distribution, four owner identities together hold \SI{25}{\percent} of
voting weight. The moderation quorum is not censorship-resistant today. &
Accepted knowingly, with growth as the expected remedy. The paper notes the tension: the network's
answer to weak demand is contraction, and contraction concentrates. Censorship resistance is, in
effect, a function of demand for the cloud. \\
\textbf{No chain-level privacy} &
The shielded pool is being retired. Every payment will be publicly linkable. &
A deliberate trade for a smaller, auditable chain. Payment privacy becomes an open problem, and the
paper says so rather than claiming otherwise. \\
\textbf{One discovery endpoint} &
The service that tells the network which applications exist and where they run is a single
Foundation-operated API, with no on-chain or peer-to-peer fallback. &
If it goes down, nothing new propagates; the last-known configuration keeps serving. A real
control-plane single point of failure, stated in the full paper's centralisation inventory. \\
\textbf{Four keys can produce blocks} &
An emergency path lets a 2-of-4 key set produce blocks that bypass the node-eligibility check, with
no time or stall-detection gating on when it may be used. &
Intended for recovery from a stalled chain, and the set is planned to move to 3-of-4 with fresh
keys. But the gating is absent at either size: nothing in consensus limits \emph{when} the path
fires, only \emph{who} can fire it. The full paper lists it in the centralisation inventory
alongside the discovery endpoint. \\
\textbf{The bridge is administered} &
Mint, release and upgrades are 3-of-4 among four people from one organisation, and the contracts
are upgradable. &
Explicitly a bring-up safety net, with immutability the stated destination once a third-party audit
and a proven operating record exist. Until the upgrade path is renounced on chain, the paper
describes the bridge as administered, not trust-minimized. \\
\textbf{Accelerators attach whole} &
A GPU is passed through as an entire device. There is no subdivision, and a request under
contention is silently dropped rather than refused. &
This is the one thing the small-model thesis of \Cref{sec:short:whatfor} needs and the current
allocation cannot do. Planned with the cloud/edge merge; not shipped. \\
\bottomrule
\end{tabular}
\end{center}

None of these is hidden in the full paper, and none should be hidden here. A reader deciding whether
to run a node, deploy an application, or build on this platform is better served by knowing them
than by discovering them.
